Courts have ordered remedial redistricting in more than 12 states since 2010, appointing special masters in Alabama, Georgia, Louisiana, Maryland, New York, North Carolina, Ohio, Pennsylvania, Texas, Virginia, and Wisconsin. These remediation proceedings share a common pathology: they are slow (typically 6--18 months from invalidation to implemented plan), expensive (\$500K--\$2M in expert and attorney fees), and produce maps whose neutrality depends on the special master's individual judgment rather than a verifiable process. We analyze the existing special-master redistricting precedent and propose three model court orders for algorithmic redistricting remedies using the \texttt{bisect} platform: (1) a mandatory order requiring the state to produce an algorithmic map; (2) a special-master-input order appointing a special master who uses \texttt{bisect} to produce a map the court can adopt; and (3) an adequacy-standard order setting algorithmic criteria that the state's new map must satisfy. We show that the SHA-256 audit chain in \texttt{bisect}'s manifest format directly satisfies the reproducibility standard implicit in \textit{League of Women Voters v. Pennsylvania}, 178 A.3d 737 (Pa. 2018), and articulated as a requirement in several special-master proceedings. Key cases analyzed include \textit{Harkenrider v. Hochul} (N.Y. 2022), \textit{League of Women Voters v. Pennsylvania} (Pa. 2018), \textit{Harper v. Hall} (N.C. 2022), and \textit{Alabama Legislative Black Caucus v. Alabama} (U.S. 2015).
